\documentclass[../../main.tex]{subfiles}

\graphicspath{{../../image/}{../../figure/}}

\begin{document}

\section{自然语言}\label{section:ch12-自然语言}

在研究了Transformer架构之后，我们将探讨如何用它来处理由词、句子和段落组成的语言数据。尽管Transformer模型最初就是针对这种模态开发的，但事实证明，Transformer模型是一类非常通用的模型，并已成为大多数输入数据类型的最先进模型。本章稍后将探讨Transformer在其他领域的应用（参见12.4节）。

包括英语在内的许多语言都由一系列词（用空格分隔）和标点符号组成，语言数据是序列数据的典型代表。我们日前将重点关注词，稍后再讨论标点符号。

\subsection{词嵌入}\label{subsection:ch12-词嵌入}

我们面临的第一个挑战是将词转换成适合用作深度神经网络输入的数字表示。一种简单的方法是定义一个固定的字典，然后引入长度与这个字典大小相等的向量以及每个词的“独热”表示，其中字典中的第$k$个词用一个向量编码，该向量在第$k$个位置为1，在其他位置为0。例如，如果“aardwolf”是这个字典中的第三个词，那么它的向量表示就是$(0,0,1,0,\cdots,0)$。

独热编码有一个明显的问题，即现实中的字典可能有几十万个条目，从而导致向量的维度非常高。此外，独热编码无法捕捉词与词之间可能存在的任何相似性或关系。这两个问题都可以通过将词映射到一个低维空间来解决，这个过程叫作词嵌入（word embedding）。在这个过程中，每个词都可以表示为一个通常只有几百个维度的空间中的密集向量。

嵌入过程可以用一个大小为$D\times K$的矩阵$\mathbf{E}$来定义，其中$D$是嵌入空间的维度，$K$是字典的维度。对于每个独热编码的输入向量$\mathbf{x}_n$，我们可以用下式计算出相应的嵌入向量：

\begin{align}
  \mathbf{v}_n = \mathbf{E}\mathbf{x}_n
  \label{align:ch12-词嵌入}
\end{align}

由于$\mathbf{x}_n$具有独热编码，因此向量$\mathbf{v}_n$可以由矩阵$\mathbf{E}$的相应列给出。

我们可以使用很多方法来从文本语料库（即大型数据集）中学习矩阵$\mathbf{E}$。在此，我们将研究一种名为word2vec的流行技术（Mikolov et al., 2013），它可以视为一种简单的双层神经网络。构建一个训练集，其中每个样本都是通过考虑文本中$M$个相邻词的“窗口”获得的，一个典型的值可能是$M=5$。样本被认为是独立的，误差函数则定义为每个样本的误差函数之和。这种方法有两个模型：连续词袋模型和“跳字模型”。在连续词袋（continuous bag of words）模型中，网络训练的目标变量是中间词，其余的上下文（context）词构成输入，因此网络被训练用于“填空”。而“跳字（skip-gram）模型”将输入和输出互换了，也就是将中心词作为输入，而将目标值作为上下文词。这两个模型如图12.11所示。

模型训练完成后，嵌入矩阵$\mathbf{E}$在连续词袋模型中由第二层权重矩阵的转置矩阵给出，而在“跳字模型”中则由第一层权重矩阵给出。在嵌入空间中，语义相关的词会被映射到附近的位置。这在意料之中，因为与不相关的词相比，相关的词更有可能与相似的上下文词一起出现。例如，“城市”和“首都”两个词作为“巴黎”或“伦敦”等目标词的上下文词出现的频率可能较高，而作为“橙色”或“多项式”的上下文词出现的频率较低。如果将“巴黎”和“伦敦”映射到附近的嵌入向量上，网络就能更容易地预测缺失词的概率。

事实证明，习得的嵌入空间往往比相关词的邻近性具有更丰富的语义结构，这使得进行简单的向量运算成为可能。例如，“巴黎之于法国就像罗马之于意大利”这一概念可以通过对嵌入向量的运算来表达。如果我们用$\mathbf{v}(\mathrm{word})$来表示“词”的嵌入向量，则可以得到

\begin{align}
  \mathbf{v}(\mathrm{Paris}) - \mathbf{v}(\mathrm{France}) + \mathbf{v}(\mathrm{Italy}) \approx \mathbf{v}(\mathrm{Rome})
  \label{align:ch12-词嵌入类比}
\end{align}

词嵌入本身最初是作为自然语言处理工具开发的。如今，它们更多地用作深度神经网络的预处理步骤。在这方面，它们可以被视为深度神经网络的第一层。它们可以被一些标准的预训练嵌入矩阵来固定，也可以作为一个自适应层来处理，作为系统整体端到端训练的一部分来学习。在后一种情况下，可以使用随机权值或标准嵌入矩阵对嵌入层进行初始化。

\subsection{分词}\label{subsection:ch12-分词}

使用固定字典的一个问题是，我们无法处理字典中没有的词或拼写错误的词。此外，这种方法不考虑标点符号或其他字符序列，如计算机代码。解决这些问题的另一种方法是在字符层面进行操作而不是使用词来操作，这样我们的字典就包含了大小写字母、数字、标点符号以及其他符号（如空格和制表符等）。但这种方法的缺点是，它忽略了语言中语义上重要的词结构，因此后续的神经网络必须学会从基本字符中重新组合词。此外，对于给定的文本，还需要执行更多的连续步骤，从而增加了处理序列的计算成本。

我们可以通过使用预处理步骤将一串词和标点符号转换成一串token，从而将字符级和词级表示的优势结合起来。通过进行这种分词，系统便能处理其他类型的序列，如计算机代码，甚至处理其他模态，如图像（参见12.4.1小节）。这也意味着同一个词的变体可以有相关的表示。例如，“cook”“cooks”“cooked”“cooking”和“cooker”都是相关的并共享共同元素“cook”，而“cook”本身也可以表示为token。

有许多方法可以实现分词。例如，用于数据压缩的字节对编码（byte pair encoding）技术可以通过连接字符而非字节来适应文本分词（Sennrich, Haddow, and Birch, 2015）。首先用单个字符列表对token列表进行初始化。然后搜索文本中出现频率最高的相邻token对，并用新token替换这些token。为确保词不被连接，如果第二个token以空格开头，则新token不是由两个token组成的。反复进行上述过程，如图12.12所示。最初，token数量等于字符数量，因而相对较少。随着token的形成，token的总数也会增加，如果持续时间足够长，token最终将与文本中的词集相对应。作为字符级和词级表示之间的折中方案，token的总数通常是预先确定的。当达到这个数量时，算法就会停止。

在将深度学习运用到自然语言的实际操作中时，输入文本通常首先映射为token化的表示。不过，在本章的其余部分，我们将使用词级表示法，因为这样更容易说明和引出关键概念。

\subsection{词袋模型}\label{subsection:ch12-词袋模型}

现在，我们要对自然语言中的词（或token）等有序向量序列的联合分布$p(x_1,\cdots,x_N)$进行建模。最简单的方法是假定词是从同一分布中独立抽取的，因此联合分布被全分解，形式为

\begin{align}
  p(x_1,\cdots,x_N) = \prod_{n=1}^{N} p(x_n)
  \label{align:ch12-词袋模型}
\end{align}

这可以表示为一个概率图模型，其中节点是孤立的，相互之间无连接。我们可以使用词袋法构建一个简单的文本分类器。例如，在情感分析中，可以将一段代表餐厅评论的文本分为正面评论和负面评论。朴素贝叶斯分类器假定每个类$C_k$中的词是独立的，但每个类的词分布不同，因此

\begin{align}
  p(x_1,\cdots,x_N\mid C_k) = \prod_{n=1}^{N} p(x_n\mid C_k)
  \label{align:ch12-词袋类条件}
\end{align}

在给定先验类概率$p(C_k)$的前提下，新序列的后验类概率为

\begin{align}
  p(C_k\mid x_1,\cdots,x_N) \propto p(C_k)\prod_{n=1}^{N} p(x_n\mid C_k)
  \label{align:ch12-词袋后验}
\end{align}

类-条件密度$p(x_n\mid C_k)$和先验概率$p(C_k)$都可以使用训练集的频率进行估计。只需要将这些概率中的每个概率设置为该词在训练集中出现的次数，就能得到该模型的最大似然解。这就是所谓的词袋（bag-of-words）模型，因为它完全忽略了词的排序。对于一个新的序列，将表格条目相乘以得到所需要的后验概率。请注意，如果测试集中出现一个训练集中没有的词，那么相应的概率估计值将为零，因此通常会在训练后对这些估计值进行“平滑”，即在所有条目中均匀地重新分配一个较小的概率水平，以避免出现零值。

\subsection{自回归模型}\label{subsection:ch12-自回归模型}

词袋模型的一个主要局限在于完全忽略了词序。为了解决这个问题，我们可以采用自回归方法。在不失通用性的前提下，我们可以将词序列的分布分解为条件分布的乘积，其形式为

\begin{align}
  p(x_1,\cdots,x_N) = p(x_1)\prod_{n=2}^{N} p(x_n\mid x_1,\cdots,x_{n-1})
  \label{align:ch12-自回归语言模型}
\end{align}

这可以用一个概率图模型（参见图11.29）来表示，其中序列里的每个节点都会从各自的前一个节点接收一个连接。我们可以用一个表格来表示式~\eqref{align:ch12-自回归语言模型}右侧的每一项，该表格中的条目也是通过训练集的简单频率计数估算出来的。然而，这些表格的大小会随着序列长度的增加呈指数增长（见习题12.12），因此这种方法的成本会高得令人望而却步。

我们可以假设式~\eqref{align:ch12-自回归语言模型}右侧的每个条件分布都与之前的所有观测值无关，但最近的$L$个词除外，从而大大简化模型。例如，如果$L=2$，那么在该模型下，一组$N$个观测值的联合分布为

\begin{align}
  p(x_1,\cdots,x_N) = p(x_1)\,p(x_2\mid x_1)\prod_{n=3}^{N} p(x_n\mid x_{n-1},x_{n-2})
  \label{align:ch12-n元模型}
\end{align}

在相应的图模型中，每个节点都有前两个节点的连接（参见图11.30）。此处，我们假设所有变量共享条件分布$p(x_n\mid x_{n-L:n-1})$。同样，式~\eqref{align:ch12-n元模型}右侧的每个分布都可以用表格来表示，其值是从抽取自训练语料的连续词的三元组统计中估算出来的。$L=1$的情况称为二元（bi-gram）模型，因为它取决于相邻的词对。同样，$L=2$的情况涉及相邻词的三元组，叫作三元（tri-gram）模型。一般来说，这些模型都叫作$n$元（$n$-gram）模型。

本节到目前为止讨论的所有模型都可以生成式运行，以合成新文本。例如，如果我们提供一个序列中的第1和第2个词，那么我们就可以从三元统计$p(x_n\mid x_{n-1},x_{n-2})$中抽样生成第3个词，然后我们可以使用第2和第3个词来抽样生成第4个词，以此类推。然而，由于每个词都是根据前两个词预测出来的，因此产生的文本将是不连贯的。高质量的文本模型必须考虑语言中的远距离依赖关系。另一方面，我们不能简单地类推。然而，由于每个词都是根据前两个词预测出来的，因此产生的文本将是不连贯的。另一方面，我们不能简单地通过增加$L$的值，因为概率表的大小会呈指数级增长，所以成本远远超过三元模型。不过，当我们考虑现代语言模型时，自回归表示法将发挥核心作用，这些模型并不基于概率表，而基于配置为Transformer的深度神经网络。

\subsection{递归神经网络}\label{subsection:ch12-递归神经网络}

$n$元模型等技术在序列长度上的缩放性很差，因为它们存储的是完全通用的条件概率表，而基于神经网络的参数化模型，我们可以实现更好的缩放。设想我们仅仅用标准的前馈神经网络来处理自然语言中的词序列。此时便会出现一个问题，即网络有固定数量的输入和输出，而我们需要网络能够处理训练集和测试集中长度可变的序列。此外，如果序列中某个位置的一个词或一组词代表某个概念，那么处在不同位置的同一个词或同一组词很可能在新的位置代表相同的概念。这会让我们想起之前处理图像数据时遇到的等变性（参见第10章）。如果我们能够构建一个能在序列中共享参数的网络架构，那么我们不仅可以捕捉到等变性，还可以大大减少模型中自由参数的数量，并处理不同长度的序列。

为了解决这个问题，我们可以借鉴隐马尔可夫模型，引入一个与序列中每一步$n$相关的显式隐变量$\mathbf{z}_n$。神经网络将当前词$\mathbf{x}_n$和当前隐状态$\mathbf{z}_{n-1}$作为输入，并产生输出词$\mathbf{y}_n$以及隐变量的下一个状态$\mathbf{z}_n$。然后，我们可以将该网络的各个副本连接起来，其中权重值在各个副本之间共享。由此产生的架构称为递归神经网络（Recurrent Neural Network, RNN），如图12.13所示。此处隐状态的初始值可以初始化为某个默认值，例如$\mathbf{z}_0=(0,0,\cdots,0)$。

图12.13显示了带有参数$\mathbf{w}$的一般递归神经网络。它将$\mathbf{x}_1,\cdots,\mathbf{x}_N$作为输入，并生成一个序列$\mathbf{y}_1,\cdots,\mathbf{y}_N$作为输出。图中的每个方框对应一个具有非线性隐藏单元的多层网络。关于如何在实践中使用递归神经网络，请参考将句子从英语翻译成荷兰语的具体任务。句子的长度是可变的，且每个输出句子的长度可能与相应的输入句子的具体任务。输入句子的长度是可变的，且网络在开始生成输出句子之前，可能需要看到整个输入句子。

输入句子。我们可以使用递归神经网络解决这个问题，方法是输入完整的英文句子，然后输入一个特殊的token（用$\langle\mathrm{start}\rangle$表示），以触发翻译过程的开始。在训练过程中，网络将学会如何将$\langle\mathrm{start}\rangle$与输出句子的开头联系起来。我们还可以将每个连续生成的词输入下一个时间步的输入中，如图12.14所示。训练网络以生成特定的$\langle\mathrm{stop}\rangle$token，以表示翻译完成。网络的前几个阶段用于吸收输入序列，而相关的输出向量则被忽略。网络的这一部分可以看作一个“编码器”，其中整个输入句子都压缩成了隐变量的状态$\mathbf{z}^*$。其余的网络阶段则充当“解码器”，旨在逐字生成翻译好的句子作为输出。请注意，每个输出词都会作为输入词输入网络的下一阶段，因此这种方法具有类似于式~\eqref{align:ch12-自回归语言模型}的自回归结构。

解决梯度消失和梯度爆炸问题以及有限远距离依赖关系的一种方法是修正神经网络架构，从而允许额外的信号路径，该路径可以绕过网络每个阶段的许多处理步骤，进而允许在更多的时间步内记忆信息。长短时记忆（Long Short-Term Memory, LSTM）模型（Hochreiter and Schmidhuber, 1997）和门控循环单元（Gated Recurrent Unit, GRU）模型（Cho et al., 2014）是最广为人知的示例。虽然与标准递归神经网络相比，它们的性能有所提高，但在远距离依赖关系建模方面，它们的能力仍然有限，而且每个单元的额外复杂性意味着长短时记忆模型的训练速度相比标准递归神经网络更慢。此外，所有递归神经网络的信号路径都与序列中的步数成线性关系。最后，由于处理过程具有连续性，它们不支持在单个训练示例内进行并行计算。特别是，这意味着递归神经网络难以有效利用基于图形处理单元的现代高度并行硬件。用Transformer取代递归神经网络可以解决这些问题。

\subsection{通过时间反向传播}\label{subsection:ch12-通过时间反向传播}

与普通神经网络一样，递归神经网络可根据反向传播计算出的梯度，通过随机梯度下降法进行训练，并通过自动微分法进行计算。误差函数包括所有输出单元对每个单元误差的求和，其中每个输出单元都有一个softmax激活函数和一个相关的交叉熵误差函数。在前向传播过程中，激活值从序列中的第一个输入节点一直传播到序列中的所有输出节点，然后误差信号沿着相同的路径反向传播。这一过程称为通过时间反向传播（backpropagation through time），其原理简单明了。然而在实践中，由于深度网络架构会出现梯度消失或梯度爆炸的问题，因此训练很长的序列可能会很困难。

标准递归神经网络的另一个问题是，其无法很好地处理远距离依赖关系。这在自然语言中尤为突出，因为这种依赖关系非常普遍。在较长的文本段落中，不妨引入一个概念，这个概念对预测文本后面出现的词起着重要作用。在图12.14所示的架构中，英文句子的整个概念必须在长度固定的单个隐向量$\mathbf{z}^*$内被捕捉到，而如果序列较长，这个问题就会越来越严重。这就是所谓的瓶颈问题（bottleneck problem），因为任意长度的序列都必须总结为单个激活的隐向量，而网络只有在处理完整个输入序列后才能开始生成输出。

\end{document}
